\chapter{Results, Evaluation and Discussion}\label{ch:results}

This chapter presents and discusses the results of the study. \Cref{ch:eda} answers the descriptive question
\emph{what happened?} through exploratory analysis. \Cref{ch:models} reports the production models for
forecasting, classification, segmentation and anomaly detection, and \cref{ch:algorithms} studies the nine
algorithms one by one. The chapter then presents the decision-support insights, evaluates the system as a
whole, decides the hypotheses, answers the research questions and compares the findings with the literature.

\section{Exploratory data analysis}\label{ch:eda}

This section answers the descriptive question, \emph{what happened?}, and begins to answer \emph{why}. All
results were computed on the full cleaned dataset by \texttt{eda.py}.

\subsection{Overall performance and growth}
Over the whole period, the company sold \TotalUnits million units in \NRows orders, worth \BDT\TotalSales
million in net sales. The average order was worth \BDT\AOV, and discounts averaged \AvgDiscount\% of gross
value, a cost of \BDT\DiscountCost million. \Cref{tab:annual} and \cref{fig:trend} show steady growth. Between
the first and the last complete year (2020--2025), net sales grew at a compound annual rate of \CAGRSales\%.
Sales in January--September 2026 were \YTDGrowth\% above the same months of 2025.

\begin{table}[htbp]
\centering
\small
\caption{Annual performance summary, 2020--2026}
\label{tab:annual}
\begin{adjustbox}{max width=\linewidth}
\begin{tabular}{lrrrrrrr}
\toprule
\textbf{Year} & \textbf{Orders} & \textbf{Units} & \textbf{Net sales (BDT M)} & \textbf{AOV (BDT)} & \textbf{Avg price} & \textbf{Disc. \%} & \textbf{YoY \%} \\
\midrule
2020 & 74,788 & 398,005 & 28.62 & 382.7 & 72.0 & 2.34 & -- \\
2021 & 73,636 & 418,104 & 37.02 & 502.7 & 88.6 & 2.35 & +29.3 \\
2022 & 73,898 & 440,150 & 41.00 & 554.8 & 93.4 & 2.34 & +10.8 \\
2023 & 73,598 & 460,465 & 46.17 & 627.3 & 100.5 & 2.36 & +12.6 \\
2024 & 74,282 & 486,293 & 53.03 & 713.9 & 109.8 & 2.35 & +14.9 \\
2025 & 74,442 & 507,262 & 61.78 & 829.9 & 122.3 & 2.35 & +16.5 \\
2026* & 55,356 & 381,678 & 50.75 & 916.8 & 133.8 & 2.34 & -- \\
\bottomrule
\end{tabular}
\end{adjustbox}
\par\smallskip\parbox{0.92\linewidth}{\footnotesize\emph{Note:} * 2026 covers January--September only, so growth is not computed for it. AOV = average order value; price = mean unit selling price.}
\end{table}

\begin{figure}[htbp]\centering
\includegraphics[width=\linewidth]{eda_monthly_trend}
\caption{Monthly net sales with 12-month moving average (top) and monthly units sold (bottom)}\label{fig:trend}
\end{figure}

\subsubsection{Where does growth come from?}
Units grew by only \CAGRUnits\% a year, whereas the average selling price grew by \CAGRPrice\% a year. Sales
growth can be decomposed exactly as $(1+g_{\text{sales}}) = (1+g_{\text{volume}})(1+g_{\text{price}})$.
\Cref{tab:pricevolume} and \cref{fig:pricevolume} show that, on average, about \AvgVolGrowth percentage points
of annual growth came from volume and \AvgPriceGrowth from price and product mix. The exceptional rise in 2021
coincides with the launch of higher-priced gift boxes, a mix effect. Later years reflect regular list-price
increases. This distinction matters for management: price-led growth cannot continue indefinitely without
eventually testing consumers' willingness to pay.

\begin{table}[htbp]
\centering
\small
\caption{Decomposition of annual sales growth into volume and price effects}
\label{tab:pricevolume}
\begin{adjustbox}{max width=\linewidth}
\begin{tabular}{lrrr}
\toprule
\textbf{Year} & \textbf{Sales growth \%} & \textbf{Volume growth \%} & \textbf{Avg price growth \%} \\
\midrule
2021 & +29.3 & +5.0 & +23.1 \\
2022 & +10.8 & +5.3 & +5.2 \\
2023 & +12.6 & +4.6 & +7.6 \\
2024 & +14.9 & +5.6 & +8.8 \\
2025 & +16.5 & +4.3 & +11.7 \\
\bottomrule
\end{tabular}
\end{adjustbox}
\par\smallskip\parbox{0.92\linewidth}{\footnotesize\emph{Note:} $(1+g_{sales}) = (1+g_{volume})(1+g_{price})$; price is net revenue per unit and therefore also reflects changes in product mix.}
\end{table}

\begin{figure}[htbp]\centering
\includegraphics[width=\linewidth]{eda_price_volume}
\caption{Decomposition of growth (left) and price versus volume indices (right)}\label{fig:pricevolume}
\end{figure}

\subsubsection{Is demand sensitive to price?}
To separate the effect of price from that of general growth, the price elasticity of demand was estimated
with a fixed-effects regression on the product-month panel \citep{wooldridge2020}:
\[
  \log(\text{units}_{pt}) = \beta \log(\text{price}_{pt}) + \alpha_p + \gamma_{\text{year}(t)} + \delta_{\text{month}(t)} + \varepsilon_{pt}.
\]
The product effects $\alpha_p$ absorb differences in size between products. The year effects absorb common
market growth, and the month effects absorb seasonality. The estimated elasticity is
$\hat\beta = \Elasticity$, with a 95\% confidence interval from \ElasticityLo to \ElasticityHi (robust standard
errors; $p = \ElasticityP$; $R^2 = \ElasticityRtwo$). It is not significantly different from zero. Over the
observed range, demand therefore appears \emph{price-inelastic}. A naive regression without year effects gives
a misleading \emph{positive} elasticity of \ElasticityNaive, because prices and volumes both rose over time.
This is a textbook illustration of confounding, and the reason explanatory models need careful
specification \citep{shmueli2010}.

\subsection{Seasonality and calendar effects}
Sales show a clear annual cycle (\cref{fig:seasonality}). \PeakMonth is the strongest month, with a seasonal
index of \PeakIndex (average month = 100). October (\OctIndex) and November (\NovIndex) are also above average,
consistent with the year-end festive and gifting season. \TroughMonth is the weakest month (\TroughIndex). Part
of February's weakness is simply its shorter length: per trading day its index is \FebPerDay, while December's
per-day index remains high at \DecPerDay. A Kruskal--Wallis test on detrended daily sales confirms that
calendar month matters ($H = \KWMonthH$, $p<0.001$). By contrast, the day of the week has almost no effect: the
coefficient of variation of average daily sales across weekdays is only \DowCV\%.

\begin{figure}[htbp]\centering
\includegraphics[width=\linewidth]{eda_seasonality}
\caption{Seasonal index by month (top) and year $\times$ month heat map of net sales (bottom)}\label{fig:seasonality}
\end{figure}

\subsection{Products and categories}
\Cref{tab:category} and \cref{fig:pareto} show the category structure. \TopCategory is the largest category
(\TopCategoryShare\% of sales), followed by Pound Cake and Swiss Roll. Layer Cake has the highest growth rate
(\FastCategoryCAGR\% a year). This growth comes mainly from gift-box variants added to the range in 2021, a
reminder that category growth figures combine organic growth with range changes. The Pareto analysis divides
the \NProducts products into \NClassA A-class products, which earn 80\% of revenue, and \NClassB B-class and
\NClassC C-class products. The ten best-selling products alone deliver \TopTenShare\% of net sales
(\cref{tab:topproducts}).

\begin{table}[htbp]
\centering
\small
\caption{Revenue and growth by product category}
\label{tab:category}
\begin{adjustbox}{max width=\linewidth}
\begin{tabular}{lrrrr}
\toprule
\textbf{Category} & \textbf{Net sales (BDT M)} & \textbf{Share \%} & \textbf{Products} & \textbf{CAGR 2020-25 \%} \\
\midrule
Muffin & 71.47 & 22.4 & 9 & +32.3 \\
Pound Cake & 61.76 & 19.4 & 12 & +9.7 \\
Swiss Roll & 38.81 & 12.2 & 4 & +9.6 \\
Fruit Cake & 28.90 & 9.1 & 5 & +9.9 \\
Layer Cake & 28.47 & 8.9 & 4 & +67.7 \\
Dry Cake & 25.59 & 8.0 & 3 & +9.4 \\
Cookie & 21.72 & 6.8 & 3 & +11.0 \\
Plain Cake & 16.57 & 5.2 & 6 & +10.4 \\
Velvet Cake & 11.75 & 3.7 & 1 & +10.3 \\
Classic Brownie & 6.27 & 2.0 & 1 & +10.8 \\
Premium & 4.02 & 1.3 & 1 & -- \\
Breakfast & 1.75 & 0.6 & 1 & -- \\
Bread & 1.28 & 0.4 & 2 & -- \\
\bottomrule
\end{tabular}
\end{adjustbox}
\par\smallskip\parbox{0.92\linewidth}{\footnotesize\emph{Note:} CAGR is computed between the first and last complete years. Categories launched after 2020 show growth from their launch year base.}
\end{table}
\begin{table}[htbp]
\centering
\small
\caption{Top ten products by net sales (Pareto/ABC analysis)}
\label{tab:topproducts}
\begin{adjustbox}{max width=\linewidth}
\begin{tabular}{lrrr}
\toprule
\textbf{Product} & \textbf{Net sales (BDT M)} & \textbf{Cumulative share \%} & \textbf{Class} \\
\midrule
MUF007 Cappuccino Muffin Gift Pack & 15.41 & 4.8 & A \\
MUF006 Vanilla Muffin Gift Box & 12.48 & 8.8 & A \\
LAY004 Vanilla Layer Cake Gift Box & 12.44 & 12.7 & A \\
LAY002 Chocolate Layer Cake Gift Box & 12.29 & 16.5 & A \\
MUF003 Chocolate Muffin Gift Box & 12.25 & 20.4 & A \\
MUF009 Chocochips Muffin & 11.99 & 24.1 & A \\
MUF008 Raisin Muffin & 11.80 & 27.8 & A \\
VEL001 Chocolate Velvet Cake & 11.75 & 31.5 & A \\
SWR002 Chocolate Swiss Roll & 11.68 & 35.2 & A \\
SWR004 Strawberry Swiss Roll & 10.44 & 38.5 & A \\
\bottomrule
\end{tabular}
\end{adjustbox}
\end{table}

\begin{figure}[htbp]\centering
\includegraphics[width=\linewidth]{eda_category_pareto}
\caption{Share of net sales by category (left) and Pareto curve of products (right)}\label{fig:pareto}
\end{figure}

\subsection{Regions, channels and customers}
Sales are geographically concentrated. Dhaka generates \DhakaShare\% of net sales, and Dhaka and Chattogram
together \DhakaCtgShare\% (\cref{tab:segments}). General Trade is the main channel (\GTShare\%), followed by
Modern Trade (\MTShare\%), Distributors (\DistShare\%) and E-commerce (\EcomShare\%). The channels differ
sharply in order size. A distributor order averages \DistQty units, an E-commerce order \EcomQty, General Trade
\GTQty and Modern Trade \MTQty (\cref{fig:regionchannel}). Customer type, on the other hand, makes no difference
to order size (Kruskal--Wallis $p=\PCustType$), and average order values are almost identical across customer
types. Order size is therefore governed by the \emph{route to market}, not by who the customer is.

\begin{table}[htbp]
\centering
\footnotesize
\caption{Performance by region, channel and customer type}
\label{tab:segments}
\begin{adjustbox}{max width=\linewidth}
\begin{tabular}{llrrrrr}
\toprule
\textbf{Dimension} & \textbf{Level} & \textbf{Net sales (BDT M)} & \textbf{Share \%} & \textbf{Orders} & \textbf{AOV (BDT)} & \textbf{Units/order} \\
\midrule
Region & Dhaka & 139.76 & 43.9 & 169,661 & 823.8 & 7.85 \\
Region & Chattogram & 45.89 & 14.4 & 70,440 & 651.4 & 6.33 \\
Region & Gazipur & 25.74 & 8.1 & 39,853 & 645.9 & 6.34 \\
Region & Narayanganj & 21.23 & 6.7 & 35,195 & 603.1 & 5.93 \\
Region & Cumilla & 15.71 & 4.9 & 29,911 & 525.2 & 5.19 \\
Region & Sylhet & 15.27 & 4.8 & 29,901 & 510.6 & 4.92 \\
Region & Rangpur & 12.23 & 3.8 & 30,012 & 407.6 & 4.16 \\
Region & Rajshahi & 12.03 & 3.8 & 25,152 & 478.2 & 4.79 \\
Region & Khulna & 11.73 & 3.7 & 25,094 & 467.3 & 4.71 \\
Region & Mymensingh & 10.76 & 3.4 & 24,849 & 432.9 & 4.34 \\
Region & Barishal & 8.04 & 2.5 & 19,932 & 403.2 & 3.96 \\
Channel & General Trade & 145.21 & 45.6 & 260,431 & 557.6 & 5.32 \\
Channel & Modern Trade & 81.55 & 25.6 & 124,619 & 654.4 & 6.28 \\
Channel & Distributor & 72.90 & 22.9 & 64,861 & 1,123.9 & 11.43 \\
Channel & E-commerce & 18.72 & 5.9 & 50,089 & 373.7 & 3.66 \\
Customer type & Retailer & 143.28 & 45.0 & 224,582 & 638.0 & 6.19 \\
Customer type & Super Shop & 63.31 & 19.9 & 100,073 & 632.7 & 6.16 \\
Customer type & Distributor & 47.91 & 15.0 & 75,413 & 635.3 & 6.19 \\
Customer type & Online Customer & 31.95 & 10.0 & 49,759 & 642.0 & 6.17 \\
Customer type & Corporate & 31.93 & 10.0 & 50,173 & 636.4 & 6.21 \\
\bottomrule
\end{tabular}
\end{adjustbox}
\end{table}

\begin{figure}[htbp]\centering
\includegraphics[width=\linewidth]{eda_region_channel}
\caption{Net sales (left) and mean units per order (right) by region and channel}\label{fig:regionchannel}
\end{figure}

\subsection{Distributions and relationships}
Order quantity and order value are strongly right-skewed, with skewness of \QtySkew and \ValueSkew
respectively (\cref{fig:distributions}). Most orders are small, and a long tail of bulk orders exists. These
shapes justify the use of medians, rank-based tests and log transformations throughout the analysis. The
Spearman correlation matrix (\cref{fig:correlation}) shows that order value is driven by both quantity
($\rho=\CorrSalesQty$) and price ($\rho=\CorrSalesPrice$). Quantity and price, however, are practically
unrelated ($\rho=\CorrQtyPrice$): customers do not buy fewer packs of the more expensive products.

\begin{figure}[htbp]\centering
\includegraphics[width=\linewidth]{eda_distributions}
\caption{Distributions of units per order, order value (log scale) and discount}\label{fig:distributions}
\end{figure}

\begin{figure}[htbp]\centering
\includegraphics[width=.62\linewidth]{eda_correlation}
\caption{Spearman correlation matrix of the main numeric variables}\label{fig:correlation}
\end{figure}

\subsection{Hypothesis tests}
\Cref{tab:tests} formalises the main exploratory findings. With half a million observations, every difference
except two is statistically significant, so effect sizes are essential for interpretation. Channel explains the
most variation in order size ($\varepsilon^2=\EpsChannel$), followed by region ($\varepsilon^2=\EpsRegion$).
The association between channel and large orders is moderate (Cram\'er's $V=\CramerV$). Discount depth and
order size are \emph{uncorrelated} ($\rho=\RhoDisc$, $p=\PDisc$). The discounts given do not appear to
stimulate larger purchases. This is a potentially valuable finding, because discounts cost the company
\BDT\DiscountCost million over the period.

\begin{table}[htbp]
\centering
\small
\caption{Statistical tests of the main exploratory hypotheses}
\label{tab:tests}
\begin{adjustbox}{max width=\linewidth}
\begin{tabular}{llrrr}
\toprule
\textbf{Hypothesis (H0: no difference)} & \textbf{Test} & \textbf{Statistic} & \textbf{p-value} & \textbf{Effect size} \\
\midrule
Units/order by channel & Kruskal-Wallis & 43,269.88 & <0.001 & 0.0865 \\
Units/order by region & Kruskal-Wallis & 27,261.36 & <0.001 & 0.0545 \\
Units/order by customer type & Kruskal-Wallis & 2.62 & 0.623 & -0.0000 \\
Daily sales by month (detrended) & Kruskal-Wallis & 510.92 & <0.001 & 0.2293 \\
Discount \% vs units/order & Spearman rho & 0.00 & 0.479 & 0.0000 \\
Large order vs channel & Chi-square & 31,777.06 & <0.001 & 0.2521 \\
\bottomrule
\end{tabular}
\end{adjustbox}
\par\smallskip\parbox{0.92\linewidth}{\footnotesize\emph{Note:} Effect sizes: $\varepsilon^2$ for Kruskal--Wallis, $\rho^2$ for Spearman and Cram\'er's $V$ for chi-square. With $n=500{,}000$ almost any difference is `significant', so the effect size decides practical importance.}
\end{table}

\subsection{Summary of exploratory findings}
\begin{enumerate}
  \item Revenue growth is strong (\CAGRSales\% a year) but mostly price- and mix-driven; volume grows at \CAGRUnits\%.
  \item Demand is price-inelastic in the observed range, and discounts show no measurable effect on order size.
  \item Seasonality is pronounced, with a \PeakMonth peak and a \TroughMonth trough; weekdays do not matter.
  \item Revenue is concentrated: Dhaka provides \DhakaShare\% of sales and ten products provide \TopTenShare\%.
  \item Order size depends on channel and region, not on customer type.
\end{enumerate}
These findings informed the modelling choices in the next section. Seasonality and history features are used
for forecasting, channel and region for classification, and price, growth and seasonality for segmentation.
\section{Predictive modelling results}\label{ch:models}

This section answers the predictive question, \emph{what is likely to happen next?} It covers the four
analytical problems in turn: demand forecasting, large-order classification, segmentation and anomaly detection.

\subsection{Monthly demand forecasting}
\subsubsection{Set-up}
The forecasting panel contains \PanelSeries products observed monthly. There are \PanelTrain training rows
(2020--2024) and \PanelTest test rows (January 2025 -- September 2026). Each model predicts
$\log(1+\text{units})$, and the prediction is converted back to units. The log scale suits the multiplicative
nature of demand and stops the largest products from dominating the fit. Seven algorithms were tuned by
randomised search over the expanding-window folds described in \cref{ch:method}. \Cref{app:hyper}
lists the selected hyper-parameters.

\subsubsection{One-step-ahead results}
\Cref{tab:forecast} reports test performance when each month is forecast with the actual history up to the
previous month. This is the operational scenario of a monthly planning cycle. Every machine-learning model
clearly beats both naive baselines. The best model, \BestReg, reaches RMSE \BestRegRMSE units, WAPE
\BestRegWAPE\% and $R^2=\BestRegRtwo$. That is an RMSE \GainVsSNaive\% lower than the seasonal-naive rule
(\SNaiveRMSE) and \GainVsNaive\% lower than the naive rule (\NaiveRMSE).

\begin{table}[htbp]
\centering
\footnotesize
\caption{Monthly product-demand forecasting: test-set performance (one-step-ahead, Jan 2025--Sep 2026)}
\label{tab:forecast}
\begin{adjustbox}{max width=\linewidth}
\begin{tabular}{lrrrrrrrr}
\toprule
\textbf{Model} & \textbf{MAE} & \textbf{MSE} & \textbf{RMSE} & \textbf{R$^2$} & \textbf{MAPE \%} & \textbf{WAPE \%} & \textbf{CV RMSE} & \textbf{Fit (s)} \\
\midrule
Seasonal naive (same month last year) & 120.0 & 23,233 & 152.4 & 0.551 & 15.38 & 14.74 & -- & 0.0 \\
Naive (last month) & 139.8 & 31,051 & 176.2 & 0.400 & 18.00 & 17.17 & -- & 0.0 \\
Support Vector Regression & 90.6 & 13,271 & 115.2 & 0.744 & 12.26 & 11.13 & 0.2151 & 24.9 \\
Linear Regression & 91.3 & 13,861 & 117.7 & 0.732 & 11.48 & 11.21 & 0.1691 & 0.2 \\
XGBoost & 93.3 & 14,294 & 119.6 & 0.724 & 12.28 & 11.46 & 0.1917 & 94.8 \\
Random Forest & 93.0 & 14,413 & 120.1 & 0.722 & 12.04 & 11.42 & 0.1918 & 142.6 \\
Gradient Boosting & 95.7 & 15,113 & 122.9 & 0.708 & 12.51 & 11.76 & 0.1889 & 124.9 \\
K-Nearest Neighbors & 97.0 & 15,865 & 126.0 & 0.694 & 12.01 & 11.91 & 0.1951 & 6.8 \\
Decision Tree & 103.6 & 17,978 & 134.1 & 0.653 & 13.70 & 12.73 & 0.2077 & 8.1 \\
\bottomrule
\end{tabular}
\end{adjustbox}
\par\smallskip\parbox{0.92\linewidth}{\footnotesize\emph{Note:} Errors in units per product-month. CV RMSE is the mean validation RMSE on the log scale over the expanding-window folds 2022, 2023 and 2024 used for tuning. Baselines need no fitting.}
\end{table}

Three further observations stand out. First, the differences between the better models are small. RMSE values
range only from \BestRegRMSE to \WorstRegRMSE units, and linear regression on log features (\LinRMSE) is
competitive with the tree ensembles (random forest \RFRMSE, XGBoost \XGBRMSE). Once good lag and seasonal
features exist, the problem is largely log-linear, and model complexity adds little. This agrees with
\citet{makridakis2018}. Second, the single decision tree is the weakest learner (\DTRMSE), as expected from its
high variance; averaging many trees in a forest removes most of this weakness \citep{breiman2001}. Third, the
cross-validated scores on the training years rank the models similarly to the test scores. The validation
design therefore gives a reliable guide to future performance. \Cref{fig:forecastcomp} compares the models and
shows that predictions of \BestReg track the actual values closely across products of very different size.

\begin{figure}[htbp]\centering
\includegraphics[width=\linewidth]{ml_forecast_comparison}
\caption{Test RMSE of all forecasting models (left; best model highlighted, baselines in grey) and actual versus predicted units for the best model (right)}\label{fig:forecastcomp}
\end{figure}

\subsubsection{Recursive multi-step results}
A harder test is to forecast all 21 test months from December 2024 \emph{without} seeing any actual data after
training. Each month's predictions are fed back in as history for the next month. This is the scenario of an
annual budget. \Cref{tab:recursive} shows that the best models degrade only modestly. \BestRec achieves a
WAPE of \BestRecWAPE\% ($R^2=\BestRecRtwo$), still far better than repeating the 2024 pattern. XGBoost degrades
the most under recursion. Errors compound when a model's own predictions become inputs, and boosted trees, which
fit the training data closely, are more exposed to this.

\begin{table}[htbp]
\centering
\small
\caption{Recursive 21-month forecast from December 2024 (no actual data used after training)}
\label{tab:recursive}
\begin{adjustbox}{max width=\linewidth}
\begin{tabular}{lrrrrr}
\toprule
\textbf{Model} & \textbf{MAE} & \textbf{RMSE} & \textbf{R$^2$} & \textbf{MAPE \%} & \textbf{WAPE \%} \\
\midrule
Random Forest & 91.0 & 117.0 & 0.736 & 11.85 & 11.18 \\
Support Vector Regression & 92.3 & 117.2 & 0.735 & 12.65 & 11.34 \\
Linear Regression & 94.1 & 121.4 & 0.715 & 11.73 & 11.56 \\
XGBoost & 104.1 & 137.5 & 0.635 & 16.17 & 12.79 \\
Seasonal naive (2024 pattern repeated) & 125.9 & 158.1 & 0.517 & 15.87 & 15.47 \\
\bottomrule
\end{tabular}
\end{adjustbox}
\end{table}

\subsubsection{Company-level forecasting and model reconciliation}
Product forecasts can be summed to a company total and compared with statistical models fitted directly to the
total series (\cref{tab:company}, \cref{fig:companyfc}). At this aggregate level, the Holt--Winters model with
additive trend and multiplicative seasonality is the most accurate, with a MAPE of only \HWMAPE\%. The summed
machine-learning forecasts are also accurate but slightly less so. The result is consistent with the literature:
smooth aggregate series favour statistical models, while the cross-learning ability of machine learning pays
off at the level of individual products. The framework therefore combines the two through \emph{top-down
reconciliation} \citep{athanasopoulos2009}. Holt--Winters sets the company total and its prediction interval,
and the machine-learning model distributes that total across products. Product plans then add up exactly to the
company plan.

\begin{table}[htbp]
\centering
\small
\caption{Company-level monthly units forecast, test period Jan 2025--Sep 2026}
\label{tab:company}
\begin{adjustbox}{max width=\linewidth}
\begin{tabular}{lrrr}
\toprule
\textbf{Model (company-total monthly units)} & \textbf{MAE} & \textbf{RMSE} & \textbf{MAPE \%} \\
\midrule
Holt-Winters (additive trend, multiplicative season) & 559 & 698 & 1.32 \\
SARIMA(1,1,1)(0,1,1)$_{12}$ on log units & 687 & 806 & 1.65 \\
Support Vector Regression one-step, summed over products & 712 & 875 & 1.68 \\
Random Forest recursive, summed over products & 1,186 & 1,386 & 2.77 \\
Seasonal naive (2024 repeated) & 2,549 & 2,817 & 6.00 \\
\bottomrule
\end{tabular}
\end{adjustbox}
\end{table}

\begin{figure}[htbp]\centering
\includegraphics[width=\linewidth]{ml_company_forecast_test}
\caption{Company monthly units: actual values and competing forecasts over the test period}\label{fig:companyfc}
\end{figure}

\subsubsection{What drives the forecasts?}
\Cref{fig:importance} explains the models. For \BestReg, permutation importance shows that recent history
(lags and rolling means) carries most of the information, followed by seasonal features. The XGBoost SHAP values
give the same picture and also highlight the December indicator and the 12-month lag, which capture the
seasonal peak. The price features have little influence. This matches the near-zero price elasticity found in
\cref{ch:eda}, so the models and the explanatory analysis tell a consistent story.

\begin{figure}[htbp]\centering
\includegraphics[width=\linewidth]{ml_feature_importance}
\caption{Permutation importance for the best forecasting model (left) and mean absolute SHAP values for XGBoost (right)}\label{fig:importance}
\end{figure}

\subsubsection{Twelve-month outlook}
The selected models were refitted on all data up to September 2026 and used to forecast October 2026 --
September 2027 (\cref{tab:future}, \cref{fig:future}). Net sales for these twelve months are forecast at
\BDT\FutSales million, with an 80\% interval from \FutSalesLo to \FutSalesHi million. That is \FutGrowth\% more
than the previous twelve months (\BDT\PrevSales million). Units are expected to change by \FutUnitsGrowth\%, and
the December peak is expected to recur. \FutTopCategory remains the largest category in the forecast. This
outlook can feed directly into the production, procurement and budgeting processes.

\begin{table}[htbp]
\centering
\footnotesize
\caption{Twelve-month forward forecast, October 2026 -- September 2027}
\label{tab:future}
\begin{adjustbox}{max width=\linewidth}
\begin{tabular}{lrrrrrrr}
\toprule
\textbf{Month} & \textbf{Units ML raw} & \textbf{Units (final)} & \textbf{80\% low} & \textbf{80\% high} & \textbf{Net sales (M)} & \textbf{80\% low } & \textbf{80\% high } \\
\midrule
Oct 2026 & 43,531 & 46,203 & 45,021 & 47,478 & 6.08 & 5.83 & 6.34 \\
Nov 2026 & 44,109 & 47,015 & 45,744 & 48,272 & 6.19 & 5.90 & 6.47 \\
Dec 2026 & 53,219 & 54,939 & 53,383 & 56,313 & 7.27 & 6.89 & 7.63 \\
Jan 2027 & 43,375 & 45,461 & 44,209 & 46,738 & 6.58 & 6.22 & 6.95 \\
Feb 2027 & 37,348 & 39,411 & 38,347 & 40,460 & 5.68 & 5.36 & 6.02 \\
Mar 2027 & 43,561 & 45,883 & 44,660 & 47,104 & 6.56 & 6.16 & 6.96 \\
Apr 2027 & 41,747 & 43,288 & 42,136 & 44,485 & 6.21 & 5.83 & 6.62 \\
May 2027 & 43,224 & 45,767 & 44,576 & 47,013 & 6.44 & 6.02 & 6.88 \\
Jun 2027 & 41,784 & 43,454 & 42,324 & 44,667 & 6.16 & 5.75 & 6.61 \\
Jul 2027 & 42,815 & 44,111 & 42,962 & 45,304 & 6.27 & 5.84 & 6.74 \\
Aug 2027 & 42,884 & 45,082 & 43,902 & 46,342 & 6.46 & 5.99 & 6.95 \\
Sep 2027 & 43,340 & 44,976 & 43,746 & 46,142 & 6.38 & 5.89 & 6.88 \\
\bottomrule
\end{tabular}
\end{adjustbox}
\par\smallskip\parbox{0.92\linewidth}{\footnotesize\emph{Note:} Units ML raw = sum of the Random Forest product forecasts before reconciliation. Final totals come from Holt-Winters models of company units and net sales (BDT million); intervals are the 10th and 90th percentiles of 2,000 simulated future paths. Product forecasts are scaled to these totals.}
\end{table}

\begin{figure}[htbp]\centering
\includegraphics[width=\linewidth]{ml_future_forecast}
\caption{Net sales outlook for October 2026 -- September 2027 with 80\% prediction interval}\label{fig:future}
\end{figure}

\subsection{Large-order classification}
The classifier predicts, from information available when an order is entered, whether the order will contain
at least 10 units. Large orders made up \ClsPosTrain\% of training orders and \ClsPosTest\% of test orders. The
models were trained on \ClsTrain orders (2020--2024) and tested on \ClsTest orders (2025--2026). The features
were channel, region, customer type, category, price, discount, pack size, product age and calendar variables,
plus smoothed target encodings of channel $\times$ region and of product.

\begin{table}[htbp]
\centering
\footnotesize
\caption{Large-order classification ($\geq$10 units): test-set performance, 2025--2026}
\label{tab:classification}
\begin{adjustbox}{max width=\linewidth}
\begin{tabular}{lrrrrrrrrr}
\toprule
\textbf{Model} & \textbf{Accuracy} & \textbf{Precision} & \textbf{Recall} & \textbf{F1} & \textbf{ROC-AUC} & \textbf{PR-AUC} & \textbf{Threshold} & \textbf{Lift@10\%} & \textbf{Time (s)} \\
\midrule
Logistic Regression & 0.661 & 0.357 & 0.628 & 0.455 & 0.712 & 0.418 & 0.62 & 2.35 & 15 \\
Support Vector Machine (linear) & 0.629 & 0.341 & 0.698 & 0.458 & 0.711 & 0.417 & 0.31 & 2.34 & 42 \\
Gradient Boosting (hist) & 0.709 & 0.387 & 0.502 & 0.437 & 0.711 & 0.418 & 0.23 & 2.34 & 35 \\
Random Forest & 0.628 & 0.341 & 0.701 & 0.459 & 0.710 & 0.418 & 0.51 & 2.33 & 252 \\
XGBoost & 0.635 & 0.343 & 0.676 & 0.455 & 0.709 & 0.416 & 0.20 & 2.34 & 60 \\
Decision Tree & 0.550 & 0.309 & 0.807 & 0.447 & 0.709 & 0.407 & 0.48 & 2.31 & 14 \\
K-Nearest Neighbors & 0.634 & 0.333 & 0.624 & 0.434 & 0.684 & 0.396 & 0.20 & 2.28 & 73 \\
\bottomrule
\end{tabular}
\end{adjustbox}
\par\smallskip\parbox{0.92\linewidth}{\footnotesize\emph{Note:} Positive rate in the test years: 22.5\%; a classifier that always predicts `not large' would reach 77.5\% accuracy, so accuracy alone is misleading. Thresholds maximise F1 on the 2024 validation year. KNN uses a 60,000-row training sample. Lift@10\% is the large-order rate among the 10\% highest-scored orders divided by the overall rate.}
\end{table}

\Cref{tab:classification} and \cref{fig:classification} show a clear and instructive result. All seven
algorithms reach almost the same ROC-AUC, between \ClsAUCmin and \ClsAUC, and \ClsBest is the best. The model
that always predicts ``not large'' would be \ClsMajority\% accurate yet useless. The tuned models instead
recover most large orders (recall about \ClsRecall for \ClsBest) at moderate precision. Ranked by score, the top
10\% of orders contain large orders \ClsLift times more often than average. The ceiling is set by the
\emph{information} in the data, not by the algorithm. XGBoost (\ClsXGBAUC) gains nothing over logistic
regression, and KNN (\ClsKNNAUC) is weakest because distances become uninformative over many one-hot features.

\begin{figure}[htbp]\centering
\includegraphics[width=\linewidth]{ml_classification}
\caption{ROC curves of all classifiers (left) and confusion matrix of the best model (right) for large-order classification}\label{fig:classification}
\end{figure}

The XGBoost gain importance (\cref{fig:clsimp}) identifies the reason. The sales channel accounts for about
\ClsChannelImp\% of the model's gain and region for about \ClsRegionImp\%. Product, price, discount and
calendar contribute very little. In the test years \LargeDist\% of distributor orders were large, against only
\LargeEcom\% of E-commerce orders. Order size within a channel and region is otherwise close to random. For
decision-making, then, the useful output is a \emph{channel-and-region-level} expectation of large orders for
logistics planning, rather than a precise order-by-order prediction.

\begin{figure}[htbp]\centering
\includegraphics[width=.65\linewidth]{ml_classification_importance}
\caption{Drivers of large orders: XGBoost gain aggregated by original variable}\label{fig:clsimp}
\end{figure}

\subsection{Segmentation}
\subsubsection{Product portfolio}
Each product was described by its log revenue, log average price, unit growth from 2023 to 2025, December
seasonality index and month-to-month volatility. Order size and channel mix were excluded because they barely
vary between products (\cref{ch:eda}) and would only add noise. The features were standardised. $K$-Means,
agglomerative (Ward) and Gaussian-mixture clustering were compared for $k=2,\dots,8$
(\cref{fig:clustering}, \cref{tab:clustercompare}). Within the practical range of three to five segments,
$K$-Means with $k=\ProdK$ gave the best silhouette (\ProdSil) and a Davies--Bouldin index of \ProdDB. The
resulting segments are summarised in \cref{tab:productclusters} and named automatically from their averages.

\begin{table}[htbp]
\centering
\small
\caption{Product portfolio segments (K-Means, k=5)}
\label{tab:productclusters}
\begin{adjustbox}{max width=\linewidth}
\begin{tabular}{lrrrrrr}
\toprule
\textbf{Segment} & \textbf{Products} & \textbf{Sales share \%} & \textbf{Avg price} & \textbf{Units/order} & \textbf{Unit growth \%/yr} & \textbf{Dec. index} \\
\midrule
Premium \& gift packs & 13 & 33.3 & 156 & 6.24 & +7.8 & 123 \\
Mature slow-growth lines & 13 & 31.9 & 130 & 6.17 & +2.3 & 120 \\
Core family packs & 11 & 28.1 & 133 & 6.17 & +4.9 & 134 \\
Single-serve value packs & 9 & 4.2 & 23 & 6.22 & +6.0 & 119 \\
Niche \& new lines & 6 & 2.5 & 24 & 6.22 & +4.0 & 134 \\
\bottomrule
\end{tabular}
\end{adjustbox}
\par\smallskip\parbox{0.92\linewidth}{\footnotesize\emph{Note:} Growth is the compound annual growth of units between 2023 and 2025. December index = average December units relative to the product's average month (100 = no seasonality).}
\end{table}

\subsubsection{Market segments}
The 220 combinations of region, channel and customer type were clustered on their revenue, average order
value, units per order, price, discount, growth and gift-pack share. $K$-Means selected $k=\MktK$ segments
(silhouette \MktSil). They separate clearly into \emph{core retail trade} (General and Modern Trade in the large
cities), \emph{bulk distribution}, small-ticket \emph{online and regional} cells, and a group of
\emph{emerging high-growth} cells (\cref{tab:marketclusters}).

\begin{table}[htbp]
\centering
\scriptsize
\caption{Market segments of region $\times$ channel $\times$ customer-type cells (K-Means, k=4)}
\label{tab:marketclusters}
\begin{adjustbox}{max width=\linewidth}
\begin{tabular}{lrrrrrll}
\toprule
\textbf{Segment} & \textbf{Cells} & \textbf{Sales share \%} & \textbf{AOV (BDT)} & \textbf{Units/order} & \textbf{Growth \%/yr} & \textbf{Main channels} & \textbf{Main regions} \\
\midrule
Core retail trade & 99 & 71.9 & 518 & 5.02 & +14.8 & General Trade, Modern Trade & Dhaka, Chattogram, Gazipur \\
Bulk distribution & 52 & 22.6 & 966 & 9.96 & +16.2 & Distributor & Chattogram, Dhaka, Gazipur \\
Small-ticket online \& regional & 50 & 4.2 & 341 & 3.43 & +16.7 & E-commerce, Modern Trade & Barishal, Khulna, Rajshahi \\
Emerging high-growth cells & 19 & 1.3 & 431 & 4.16 & +33.2 & E-commerce, Modern Trade & Cumilla, Khulna, Mymensingh \\
\bottomrule
\end{tabular}
\end{adjustbox}
\end{table}
\begin{table}[htbp]
\centering
\small
\caption{Comparison of clustering algorithms at the selected number of clusters}
\label{tab:clustercompare}
\begin{adjustbox}{max width=\linewidth}
\begin{tabular}{llrrrr}
\toprule
\textbf{Problem} & \textbf{Algorithm} & \textbf{k} & \textbf{Silhouette} & \textbf{Davies-Bouldin} & \textbf{Calinski-Harabasz} \\
\midrule
Products & K-Means & 5 & 0.237 & 1.271 & 16.2 \\
Products & Agglomerative (Ward) & 5 & 0.212 & 1.195 & 14.9 \\
Products & Gaussian Mixture & 5 & 0.233 & 1.237 & 15.3 \\
Market cells & K-Means & 4 & 0.243 & 1.497 & 54.2 \\
Market cells & Agglomerative (Ward) & 4 & 0.250 & 1.467 & 48.9 \\
Market cells & Gaussian Mixture & 4 & 0.165 & 1.753 & 42.4 \\
\bottomrule
\end{tabular}
\end{adjustbox}
\par\smallskip\parbox{0.92\linewidth}{\footnotesize\emph{Note:} Higher silhouette and Calinski--Harabasz and lower Davies--Bouldin values indicate better-separated clusters.}
\end{table}

\begin{figure}[htbp]\centering
\includegraphics[width=\linewidth]{ml_clustering}
\caption{Silhouette score by number of clusters (left); products (centre) and market cells (right) projected onto the first two principal components}\label{fig:clustering}
\end{figure}

The silhouette values (\ProdSil and \MktSil) indicate weak to moderate structure \citep{rousseeuw1987}. The
segments are therefore useful summaries for strategy, not sharply separated natural groups. $K$-Means and Ward
clustering agreed closely, while Gaussian mixtures performed worst on these small samples. This ranking is
reported so that a reader can judge the robustness of the segmentation.

\subsection{Anomaly detection}
\subsubsection{Daily monitoring}
For each day, an expected value was computed as a 31-day centred rolling median multiplied by the month-of-year
seasonal factor. Deviations were scored on the log scale with the robust $z$-score
$z = 0.6745\,(r-\tilde r)/\text{MAD}$ and, separately, with an Isolation Forest over four daily features.
Together the two methods flagged \AnomDays of the 2,465 trading days: \AnomDaysZ by the robust $z$-score and
\AnomDaysIF by the Isolation Forest (\cref{fig:anomdaily}, \cref{tab:dailyanom}). Most flagged days are
unusually weak trading days. In a live deployment each would be checked against known causes such as public
holidays, strikes or system outages.

\begin{figure}[htbp]\centering
\includegraphics[width=\linewidth]{anomaly_daily}
\caption{Daily net sales, expected level and detected anomalies}\label{fig:anomdaily}
\end{figure}
\begin{table}[htbp]
\centering
\small
\caption{Ten most unusual trading days (seasonally adjusted)}
\label{tab:dailyanom}
\begin{adjustbox}{max width=\linewidth}
\begin{tabular}{lrrrrl}
\toprule
\textbf{Date} & \textbf{Net sales (BDT)} & \textbf{Expected (BDT)} & \textbf{Deviation \%} & \textbf{Robust z} & \textbf{Flagged by} \\
\midrule
28 Oct 2022 & 70,215 & 113,462 & -38.1 & -4.01 & robust z \\
22 Oct 2021 & 69,194 & 109,276 & -36.7 & -3.82 & robust z \\
06 Sep 2026 & 125,262 & 195,913 & -36.1 & -3.74 & robust z \\
30 Sep 2021 & 67,358 & 103,385 & -34.8 & -3.58 & both \\
02 Feb 2022 & 69,149 & 104,763 & -34.0 & -3.47 & Isolation Forest \\
15 Aug 2021 & 66,039 & 99,664 & -33.7 & -3.44 & Isolation Forest \\
11 Feb 2025 & 106,843 & 160,093 & -33.3 & -3.38 & Isolation Forest \\
13 Jan 2023 & 84,814 & 126,483 & -32.9 & -3.34 & Isolation Forest \\
08 Jan 2025 & 113,188 & 168,043 & -32.6 & -3.31 & Isolation Forest \\
13 Mar 2025 & 117,699 & 171,855 & -31.5 & -3.17 & Isolation Forest \\
\bottomrule
\end{tabular}
\end{adjustbox}
\end{table}

\subsubsection{Transaction screening}
Without labelled errors, a detector's ability cannot be measured on the real data. Following the approach of
the preprocessing scripts, 1,000 realistic errors of four types (250 each) were therefore planted in a copy of
the 500,000 orders (\cref{tab:anomaly}). The per-product IQR rule on quantity detected only \AnomIQRRecall\% of
the planted errors, because three of the four error types do not involve quantity. Simple business rules on
price and discount caught \AnomRuleRecall\% with perfect precision but missed every quantity error. The
Isolation Forest caught \AnomIFRecall\% across all four types. Price and discount errors were found completely,
and most inflated quantities were found too. The lesson is that rules and machine learning are complementary.
Rules are precise for known failure modes, and the Isolation Forest provides broad coverage for unknown ones.

\begin{table}[htbp]
\centering
\footnotesize
\caption{Detection of 1,000 injected transaction errors in 500,000 orders}
\label{tab:anomaly}
\begin{adjustbox}{max width=\linewidth}
\begin{tabular}{lrrrrrrr}
\toprule
\textbf{Method} & \textbf{Flagged} & \textbf{Precision} & \textbf{Recall} & \textbf{Price x10} & \textbf{Qty x20} & \textbf{Disc. 30-45\%} & \textbf{Price 10\%} \\
\midrule
IQR on log quantity (per product) & 1,500 & 0.112 & 0.168 & 0.00 & 0.66 & 0.00 & 0.01 \\
Business rules (price, discount) & 750 & 1.000 & 0.750 & 1.00 & 0.00 & 1.00 & 1.00 \\
Isolation Forest (4 features) & 2,500 & 0.364 & 0.909 & 1.00 & 0.64 & 1.00 & 1.00 \\
Rules OR Isolation Forest & 2,500 & 0.364 & 0.909 & 1.00 & 0.64 & 1.00 & 1.00 \\
\bottomrule
\end{tabular}
\end{adjustbox}
\par\smallskip\parbox{0.92\linewidth}{\footnotesize\emph{Note:} The last four columns give the share of each error type that was detected. 250 errors of each type were planted in a copy of the data.}
\end{table}

\subsection{Summary of model evaluation}
\begin{table}[htbp]\centering\small
\caption{Summary of the selected model for each analytical problem}\label{tab:modelsummary}
\begin{tabularx}{\linewidth}{@{}lXl@{}}
\toprule
\textbf{Problem} & \textbf{Selected approach} & \textbf{Key test result} \\ \midrule
Product demand (1-step) & \BestReg & WAPE \BestRegWAPE\%, $R^2$ \BestRegRtwo \\
Product demand (multi-step) & \BestRec, recursive & WAPE \BestRecWAPE\% \\
Company total & Holt--Winters + top-down reconciliation & MAPE \HWMAPE\% \\
Large orders & \ClsBest & ROC-AUC \ClsAUC, lift \ClsLift \\
Product segments & $K$-Means, $k=\ProdK$ & silhouette \ProdSil \\
Market segments & $K$-Means, $k=\MktK$ & silhouette \MktSil \\
Transaction anomalies & rules OR Isolation Forest & recall \AnomComboRecall\% \\ \bottomrule
\end{tabularx}
\end{table}
As the proposal required, models were selected on appropriate metrics \emph{and} business relevance, not on
accuracy alone. Where a simpler model performed as well as a complex one, as with logistic regression versus
XGBoost for classification, the simpler and more interpretable model is preferred.

The hyper-parameters selected by the expanding-window search are listed in \cref{app:hyper}.

\section{Algorithm-level study of nine machine-learning methods}\label{ch:algorithms}

Section~7.4 of the project proposal named nine candidate algorithms: linear regression, logistic regression,
decision tree, random forest, gradient boosting, XGBoost, support vector machine, $k$-nearest neighbours and
$K$-means clustering. \Cref{ch:models} used them inside the production pipeline, with rich feature sets,
randomised hyper-parameter search and statistical benchmarks. This section takes a complementary view. Each
algorithm is studied \emph{on its own}, in a separate and fully transparent experiment, in order to show
\emph{how} it learns, which settings control it, how it fails and what it contributes to the business problem.
The experiments also serve as an independent check of \cref{ch:models}: they use a deliberately simpler feature
set and separate code, so agreement between the two studies strengthens the conclusions.

\subsection{Common experimental protocol}\label{sec:algoprotocol}
\subsubsection{Three problems}
All nine experiments use the same data and the same rules, implemented once in
\path{algorithms/common.py}:
\begin{itemize}
  \item \textbf{Problem A -- regression.} Monthly units of each product. Transactions are aggregated into one row
        per product and month, described by nine features that are known before the month starts: month, a
        December indicator, product age, mean log price, mean discount, pack weight, and three history features
        (units last month, units in the same month last year, and the mean of the previous three months).
  \item \textbf{Problem B -- classification.} Whether an order is large ($\geq 10$ units), predicted from the 18
        order-level features selected in the feature-engineering step (calendar, price, pack, history and the
        smoothed channel~$\times$~region order size). Large orders form \LargeShareTrain of training orders and
        \LargeShareTest of test orders, so the classes are imbalanced. A fixed random sample of 15,000 orders
        from 2020--2024 and 15,000 from 2025--2026 is used by every algorithm. The sample keeps the slower
        algorithms (SVM, KNN) practical and makes the comparison fair.
  \item \textbf{Problem C -- clustering.} One profile per product (\KMProducts products) from 2020--2024: average
        price, units per month on sale, average order size, average discount, share of large orders and the
        share of December sales. There is no target variable.
\end{itemize}

\subsubsection{Four-way time split}
Every supervised experiment follows the same sequence (\cref{fig:split}). Settings are chosen on 2024
(validation) after learning from 2020--2023. The chosen setting is then refitted on 2020--2024, and the test
period, January 2025 -- September 2026, is scored once. The test data therefore never influence any choice.
Every model is compared with a naive baseline: ``units this month = units last month'' for problem~A and
``always small'' for problem~B.

\begin{figure}[htbp]\centering
\begin{tikzpicture}[x=1.75cm, every node/.style={font=\footnotesize}]
\foreach \y/\lab in {0/2020,1/2021,2/2022,3/2023,4/2024,5/2025,6/2026} \node at (\y+0.5,0.85) {\lab};
\fill[brand!25] (0,0) rectangle (4,0.5);   \node at (2,0.25) {\textbf{fit} -- learn while tuning};
\fill[accent!35] (4,0) rectangle (5,0.5);  \node at (4.5,0.25) {\textbf{val}};
\fill[brand!45] (0,-0.75) rectangle (5,-0.25); \node at (2.5,-0.5) {\textbf{train} -- learn again with the chosen setting};
\fill[green!30!gray!40] (5,-0.75) rectangle (6.75,-0.25); \node at (5.875,-0.5) {\textbf{test}};
\draw[thick,-{Stealth}] (0,-1.05) -- (7,-1.05) node[right] {time};
\node[align=left, anchor=west, font=\scriptsize] at (0,-1.45) {1. choose settings on 2024 \quad 2. refit on 2020--2024 \quad 3. score once on 2025--2026};
\end{tikzpicture}
\caption{Time-ordered protocol used by every supervised experiment in this section}\label{fig:split}
\end{figure}

\subsubsection{Implementation}
Each algorithm has its own script in the folder \path{algorithms/} (\cref{tab:lessons}). Every script is divided
into small numbered steps, so it can be run as a whole or cell by cell. All steps print their results and save
their charts to \path{algorithms/outputs/}. The final test scores and the key numbers quoted below are written
to JSON files, and \path{algorithms/export_to_report.py} turns them into the tables and \LaTeX{} macros of this
section. As in the rest of the report, no number here is typed by hand.

\begin{table}[htbp]\centering\small
\caption{Algorithm experiments and what each one demonstrates}\label{tab:lessons}
\begin{tabularx}{\linewidth}{@{}llX@{}}
\toprule
\textbf{Script} & \textbf{Problem} & \textbf{Concepts demonstrated} \\ \midrule
\texttt{01\_linear\_regression} & A & learned equation, standardised coefficients, residual analysis \\
\texttt{02\_logistic\_regression} & B & sigmoid, odds ratios, class weights, threshold choice, ROC curve \\
\texttt{03\_decision\_tree} & B (+A) & readable rules, Gini impurity, depth and overfitting \\
\texttt{04\_random\_forest} & A (+B) & bagging, number of trees, out-of-bag score, permutation importance \\
\texttt{05\_gradient\_boosting} & A (+B) & boosting by hand, learning rate versus number of trees \\
\texttt{06\_xgboost} & B (+A) & regularisation, early stopping, rare-class weighting, gain importance \\
\texttt{07\_svm} & B (+A) & scaling, linear and RBF kernels, the $C$ and $\gamma$ settings, SVR \\
\texttt{08\_knn} & B (+A) & distance, scaling, choice of $k$, inspecting neighbours \\
\texttt{09\_kmeans\_clustering} & C & Lloyd's algorithm by hand, elbow and silhouette, naming clusters \\ \bottomrule
\end{tabularx}
\end{table}

% ============================================================================
\subsection{Linear regression}
\paragraph{Principle.} Linear regression models the target as a weighted sum of the inputs,
$\hat y = \beta_0 + \beta_1 x_1 + \dots + \beta_p x_p$. Ordinary least squares chooses the weights that minimise
the sum of squared errors, which has the closed-form solution
\[
  \hat{\boldsymbol\beta} = (\mathbf X^\top \mathbf X)^{-1}\mathbf X^\top \mathbf y .
\]
Each coefficient is the expected change in the target for a one-unit change in its input, holding the others
constant \citep{montgomery2021, james2021}.

\paragraph{Experiment.} With last month's units as the only input, the learned line is
$\widehat{\text{units}} = \LinBzero + \LinBone \times \text{units}_{t-1}$ (\cref{fig:alglin}, left). The slope
below one is \emph{regression to the mean}: an unusually strong month tends to be followed by a weaker one, and
vice versa. This single-input model already reduces the test MAE from \AlgRegBaseMAE (baseline) to \LinOneMAE
units. With all nine standardised inputs, the test MAE falls to \AlgLinMAE units ($R^2 = \AlgLinRtwo$). The
largest standardised coefficient belongs to \texttt{\LinTopFeature} (\LinTopCoef units per standard
deviation): the recent level of demand is the strongest predictor. The residual plot (\cref{fig:alglin}, right)
is centred on zero with no strong pattern, so the linear form is adequate for this problem.

\begin{figure}[htbp]\centering
\includegraphics[width=.40\linewidth]{algo/01_simple_line}\hfill
\includegraphics[width=.585\linewidth]{algo/01_residuals}
\caption{Linear regression: the learned line with one input (left); predicted versus actual units and the error distribution on the test period (right)}\label{fig:alglin}
\end{figure}

\paragraph{Assessment.} Linear regression is fast, fully explainable and, with good lag features, nearly as
accurate as the ensembles. Its weaknesses are that it cannot represent interactions or thresholds unless they
are added as features, and that correlated inputs (here the three history features) make individual
coefficients unstable even when predictions are good.

% ============================================================================
\subsection{Logistic regression}
\paragraph{Principle.} Logistic regression passes the linear score $z = \beta_0 + \boldsymbol\beta^\top\mathbf x$
through the sigmoid function to obtain a probability,
\[
  P(y=1 \mid \mathbf x) = \sigma(z) = \frac{1}{1+e^{-z}}, \qquad
  \log\frac{P}{1-P} = \beta_0 + \boldsymbol\beta^\top\mathbf x ,
\]
and the weights maximise the log-likelihood of the observed classes \citep{hosmer2013}. The exponentiated
coefficient $e^{\beta_j}$ is an \emph{odds ratio}: the factor by which the odds of a large order are multiplied
when $x_j$ increases by one unit.

\paragraph{Experiment.} With only the smoothed channel~$\times$~region order size as input, the coefficient is
\LogBone. The predicted probability of a large order rises from \LogPfive in a channel and region whose orders
average 5 units, to \LogPten at 10 units and \LogPfifteen at 15 units. With all 18 standardised inputs and
\emph{balanced class weights} (errors on the rare class count more), the same feature dominates with an odds
ratio of \LogTopOR per standard deviation. All other odds ratios are close to one. The decision threshold was
chosen by maximising F1 on 2024, which selected \LogThreshold. On the test period the model reaches ROC-AUC
\AlgLogAUC, recall \AlgLogRecall and precision \AlgLogPrec. The baseline that always answers ``small'' is
\AlgClsBaseAcc\% accurate but finds no large order at all. This shows why accuracy alone is misleading for
imbalanced classes.

\begin{figure}[htbp]\centering
\includegraphics[width=.38\linewidth]{algo/02_sigmoid}\hfill
\includegraphics[width=.60\linewidth]{algo/02_confusion_roc}
\caption{Logistic regression: the sigmoid link (left); confusion matrix and ROC curve on the test period (right)}\label{fig:alglog}
\end{figure}

\paragraph{Assessment.} Logistic regression gives calibrated probabilities and explanations that managers can
read directly (``this channel doubles the odds''). In this study it is also the most accurate classifier, which
is strong evidence that the signal is essentially one main effect.

% ============================================================================
\subsection{Decision tree}
\paragraph{Principle.} A decision tree splits the data recursively with yes/no questions on single features. At
each node, CART chooses the question that most reduces impurity \citep{breiman1984}. For classification the
Gini impurity of a node with class shares $p_k$ is $G = 1 - \sum_k p_k^2$, which is zero for a pure node. For
regression the variance (squared error) is used, and each leaf predicts the mean of its training rows.

\paragraph{Experiment.} A tree of depth two already reads like a business rule: \emph{if the channel and region
usually place orders above about six units, expect a large order}. The depth was tuned on 2024
(\cref{fig:algdt}, top). As the tree deepens, its ROC-AUC on the data it learned from keeps rising (to
\DTFitAUCDeep at depth 15), while the score on unseen 2024 orders peaks early and then falls (\DTValAUCDeep at
depth 15). This is the textbook picture of \emph{overfitting}. The best depth was \DTBestDepth, with a test
ROC-AUC of \AlgDTAUC. In the final tree \texttt{channel\_region\_avg\_qty} accounts for \DTTopShare of the
total impurity reduction. As a regressor for problem~A, a depth-5 tree reaches an MAE of \AlgDTMAE units: better
than the baseline, but the weakest of the supervised regressors.

\begin{figure}[htbp]\centering
\includegraphics[width=.63\linewidth]{algo/03_depth_tuning}\par\medskip
\includegraphics[width=\linewidth]{algo/03_tree}
\caption{Decision tree: tuning the depth on 2024 (top); the top levels of the final tree (bottom)}\label{fig:algdt}
\end{figure}

\paragraph{Assessment.} Trees are the most explainable non-linear model and need no scaling, but a single tree has
high variance: small changes in the data can produce a different tree. The two ensemble methods that follow
are designed to remove this weakness.

% ============================================================================
\subsection{Random forest}
\paragraph{Principle.} A random forest averages $B$ trees, each grown on a bootstrap sample of the rows and
choosing every split from a random subset of the features \citep{breiman2001}. If single trees have variance
$\sigma^2$ and pairwise correlation $\rho$, the variance of their average is
\[
  \operatorname{Var}\Bigl(\tfrac{1}{B}\textstyle\sum_{b=1}^{B} T_b\Bigr) = \rho\,\sigma^2 + \frac{1-\rho}{B}\,\sigma^2 ,
\]
so adding trees removes the second term, and random feature selection reduces $\rho$ \citep{hastie2009}.

\paragraph{Experiment.} A single full-depth tree has a validation MAE of \RFOneTreeMAE units on 2024. A forest
of one tree is similar (\RFValMAEone), but the error falls quickly as trees are added and levels off at about
100 trees (\RFValMAEhundred; \cref{fig:algrf}). More trees never cause overfitting; they only cost time. Each
tree leaves out about a third of the rows, and scoring each row with the trees that did not see it gives a free
\emph{out-of-bag} estimate ($R^2 = \RFOOB$). On the test period the 300-tree forest reaches an MAE of
\AlgRFMAE units ($R^2 = \AlgRFRtwo$). Permutation importance shows that shuffling
\texttt{\RFTopFeature} increases the error by \RFTopPerm units, more than any other input. The individual trees
also give a range for each forecast, which is useful for safety stock. As a classifier, the forest reaches
ROC-AUC \AlgRFAUC.

\begin{figure}[htbp]\centering
\includegraphics[width=.70\linewidth]{algo/04_n_trees}
\caption{Random forest: validation error (2024) against the number of trees}\label{fig:algrf}
\end{figure}

% ============================================================================
\subsection{Gradient boosting}
\paragraph{Principle.} Boosting builds trees \emph{in sequence}. Starting from a constant $F_0$, each new tree
$h_m$ is fitted to the negative gradient of the loss, which for squared error is simply the current residual,
and a fraction $\nu$ (the learning rate) of it is added \citep{friedman2001}:
\[
  F_m(\mathbf x) = F_{m-1}(\mathbf x) + \nu\, h_m(\mathbf x), \qquad
  h_m \approx \arg\min_h \textstyle\sum_i \bigl(y_i - F_{m-1}(\mathbf x_i) - h(\mathbf x_i)\bigr)^2 .
\]

\paragraph{Experiment.} The experiment first implements boosting by hand in six lines of code. Starting from the
mean, the training MAE falls from \GBManualStart to \GBManualFive units after only five depth-3 trees. The
library version was then tuned over the learning rate and the number of trees using staged predictions on 2024
(\cref{fig:alggb}). A high learning rate reaches its best error quickly and then overfits. A low rate needs many
more trees but degrades more slowly. The chosen setting was a learning rate of \GBRate with \GBTrees trees. The
test MAE is \AlgGBMAE units ($R^2 = \AlgGBRtwo$), and the most important input is
\texttt{\GBTopFeature}. Because the classifier version has no class-weight option, its decision threshold was
lowered to the training share of large orders. It then reaches ROC-AUC \AlgGBAUC.

\begin{figure}[htbp]\centering
\includegraphics[width=.75\linewidth]{algo/05_learning_rate}
\caption{Gradient boosting: validation error (2024) against the number of trees for three learning rates}\label{fig:alggb}
\end{figure}

% ============================================================================
\subsection{XGBoost}
\paragraph{Principle.} XGBoost is a regularised, second-order implementation of gradient boosting
\citep{chen2016}. At step $m$ it minimises
\[
  \mathcal L^{(m)} = \sum_i \ell\bigl(y_i, \hat y_i^{(m-1)} + f_m(\mathbf x_i)\bigr) + \Omega(f_m), \qquad
  \Omega(f) = \gamma T + \tfrac12 \lambda \lVert \mathbf w\rVert^2 ,
\]
where $T$ is the number of leaves and $\mathbf w$ the leaf values. The penalty $\Omega$ discourages complex trees.
Row and column sub-sampling add randomness, as in a random forest, and \emph{early stopping} ends training as
soon as the validation score stops improving.

\paragraph{Experiment.} For problem B, the weight of the rare class was set to the ratio of small to large orders
(\XGBPosWeight). With early stopping monitored on 2024, the validation ROC-AUC peaked after only \XGBBestN trees
(\XGBValAUC; \cref{fig:algxgb}). The training curve kept rising after that point, the familiar sign of
memorisation. Refitted with that many trees, XGBoost reaches a test ROC-AUC of \AlgXGBAUC and an F1 of
\AlgXGBFone. Gain importance attributes \XGBTopGain of the model's gain to the channel~$\times$~region feature.
For problem~A early stopping selected \XGBRegTrees trees, and the regressor reaches a test MAE of \AlgXGBMAE
units ($R^2 = \AlgXGBRtwo$).

\begin{figure}[htbp]\centering
\includegraphics[width=.70\linewidth]{algo/06_early_stopping}
\caption{XGBoost: ROC-AUC on the training and validation years as trees are added, with the early-stopping point}\label{fig:algxgb}
\end{figure}

% ============================================================================
\subsection{Support vector machine}
\paragraph{Principle.} A support vector machine finds the boundary with the widest margin between the classes.
The soft-margin formulation allows some violations, controlled by $C$ \citep{cortes1995}:
\[
  \min_{\mathbf w, b, \boldsymbol\xi}\ \tfrac12 \lVert\mathbf w\rVert^2 + C\sum_i \xi_i
  \quad\text{s.t.}\quad y_i(\mathbf w^\top\phi(\mathbf x_i) + b) \geq 1 - \xi_i,\ \ \xi_i \geq 0 .
\]
The kernel trick replaces inner products by a kernel such as the radial basis function
$K(\mathbf x, \mathbf x') = \exp(-\gamma\lVert\mathbf x - \mathbf x'\rVert^2)$, which lets the boundary curve
\citep{scholkopf2002}. Only the training points on or inside the margin, the \emph{support vectors}, define the
solution. Support vector regression (SVR) applies the same idea with an $\varepsilon$-insensitive tube around the
regression function.

\paragraph{Experiment.} Because the SVM works with distances, feature scaling is essential. Without it, the
validation ROC-AUC was \SVMRawAUC, barely better than guessing. With standardisation it was \SVMScaledAUC. A
linear kernel (\SVMLinAUC) was slightly better than the default RBF kernel (\SVMRbfAUC). A small grid then
selected $C = \SVMBestC$ and $\gamma = \SVMBestGamma$ (validation ROC-AUC \SVMBestValAUC;
\cref{fig:algsvm}). The small $C$ means a wide, smooth margin works best. The final model, trained on
\SVMTrainRows orders because training time grows roughly with the square of the number of rows, reaches a
test ROC-AUC of \AlgSVMAUC. It keeps \SVMNSupport of its training rows as support vectors, a sign that the two
classes overlap heavily. For problem~A the \SVRBest is better and reaches an MAE of \AlgSVMMAE units
($R^2 = \AlgSVMRtwo$).

\begin{figure}[htbp]\centering
\includegraphics[width=.70\linewidth]{algo/07_svm_tuning}
\caption{Support vector machine: validation ROC-AUC (2024) for combinations of $C$ and $\gamma$}\label{fig:algsvm}
\end{figure}

% ============================================================================
\subsection{\texorpdfstring{$k$}{k}-nearest neighbours}
\paragraph{Principle.} KNN stores the training data and predicts for a new case from its $k$ closest neighbours
$N_k(\mathbf x)$ in feature space: by majority vote for classification and by the (optionally distance-weighted)
mean for regression \citep{cover1967}:
\[
  \hat y(\mathbf x) = \frac{1}{k}\sum_{i \in N_k(\mathbf x)} y_i .
\]
There is no training step in the usual sense, but every prediction requires distances to all stored rows.

\paragraph{Experiment.} Scaling matters as much as for the SVM. Without it, the validation ROC-AUC was \KNNRawAUC,
because the distance was almost entirely determined by \texttt{units\_lag\_1}, whose values are in the
thousands. With standardisation it rose to \KNNScaledAUC. The choice of $k$ shows the bias--variance trade-off
directly (\cref{fig:algknn}). With $k=1$ the model is perfect on its own data (ROC-AUC \KNNOneFitAUC) but weak on
2024 (\KNNOneValAUC). The validation score kept improving up to the largest value tried, $k=\KNNBestK$: when the
outcome of a single order is mostly noise, averaging over many neighbours helps. The test ROC-AUC is
\AlgKNNAUC. Inspecting the neighbours of a single order shows the explanation KNN offers: ``orders like this one
were large in \emph{x} out of $k$ cases''. For problem~A, $k=\KNNRegK$ distance-weighted neighbours give an MAE
of \AlgKNNMAE units.

\begin{figure}[htbp]\centering
\includegraphics[width=.70\linewidth]{algo/08_choose_k}
\caption{$k$-nearest neighbours: ROC-AUC on the learned data and on 2024 for different $k$}\label{fig:algknn}
\end{figure}

% ============================================================================
\subsection{\texorpdfstring{$K$}{K}-means clustering}
\paragraph{Principle.} $K$-means partitions $n$ observations into $K$ clusters with centres $\boldsymbol\mu_j$ by
minimising the within-cluster sum of squares \citep{macqueen1967}:
\[
  \min_{C_1,\dots,C_K} \sum_{j=1}^{K}\ \sum_{\mathbf x_i \in C_j} \lVert \mathbf x_i - \boldsymbol\mu_j\rVert^2 .
\]
Lloyd's algorithm alternates two steps until nothing changes: assign each point to its nearest centre, then move
each centre to the mean of its points. Because the result depends on the starting centres, the algorithm is
restarted several times and the best solution is kept.

\paragraph{Experiment.} The experiment first runs Lloyd's algorithm by hand for $K=3$. It converges in \KMIters
iterations, and the within-cluster sum of squares falls at every step. $K$ was then chosen with the elbow plot
and the silhouette coefficient \citep{rousseeuw1987, kaufman1990}. The silhouette selected $K=\KMBestK$ with a
value of \KMSil, which indicates clear structure. The small cluster contains the \KMSmall bread and breakfast
products, which sell fewer units per month, in larger orders and with a higher share of large orders than the
cakes. For $K \geq 3$ the silhouette drops to about \KMSilFour, so the cakes do not form well-separated
sub-groups on these six features. This complements \cref{ch:models}, where a different feature set (revenue,
growth, seasonality and volatility) and a practical range of three to five segments produced a more detailed,
strategy-oriented segmentation. The fitted model can place a new product into a cluster from its expected
profile, so planners can treat it like similar existing products.

\begin{figure}[htbp]\centering
\includegraphics[width=.57\linewidth]{algo/09_choose_k}\hfill
\includegraphics[width=.42\linewidth]{algo/09_clusters}
\caption{$K$-means: elbow and silhouette for $K = 2,\dots,8$ (left); products and cluster centres on the first two principal components (right)}\label{fig:algkm}
\end{figure}

% ============================================================================
\subsection{Comparative results}
\Cref{tab:algoreg,tab:algocls} bring the supervised experiments together.

\begin{table}[htbp]\centering\small
\caption{Problem A, monthly units per product: test results 2025--2026 of the algorithm lessons (sorted by MAE; best in bold)}\label{tab:algoreg}
\begin{tabular}{@{}lrrrr@{}}
\toprule
\textbf{Algorithm} & \textbf{MAE} & \textbf{RMSE} & \textbf{MAPE (\%)} & \textbf{$R^2$} \\ \midrule
\textbf{XGBoost} & \textbf{96.9} & \textbf{124.9} & \textbf{13.0} & \textbf{0.699} \\
Random Forest & 97.6 & 125.4 & 12.5 & 0.696 \\
SVM & 99.8 & 125.0 & 14.0 & 0.698 \\
Linear Regression & 100.8 & 126.3 & 14.2 & 0.692 \\
Gradient Boosting & 102.6 & 132.1 & 13.9 & 0.663 \\
KNN & 102.6 & 132.1 & 13.9 & 0.663 \\
Decision Tree & 108.0 & 139.8 & 13.5 & 0.623 \\
\midrule
\emph{Baseline: same as last month} & 139.8 & 176.2 & 18.0 & 0.400 \\
\bottomrule
\end{tabular}
\end{table}

\begin{table}[htbp]\centering\small
\caption{Problem B, large order ($\geq 10$ units): test results 2025--2026 of the algorithm lessons (sorted by ROC-AUC; best in bold)}\label{tab:algocls}
\begin{tabular}{@{}lrrrrr@{}}
\toprule
\textbf{Algorithm} & \textbf{accuracy} & \textbf{precision} & \textbf{recall} & \textbf{F1} & \textbf{ROC-AUC} \\ \midrule
\textbf{Logistic Regression} & \textbf{0.592} & \textbf{0.328} & \textbf{0.765} & \textbf{0.459} & \textbf{0.718} \\
XGBoost & 0.632 & 0.348 & 0.719 & 0.469 & 0.717 \\
Decision Tree & 0.635 & 0.349 & 0.711 & 0.468 & 0.712 \\
Random Forest & 0.632 & 0.345 & 0.696 & 0.461 & 0.712 \\
SVM & 0.715 & 0.396 & 0.494 & 0.439 & 0.712 \\
Gradient Boosting & 0.557 & 0.310 & 0.788 & 0.445 & 0.693 \\
KNN & 0.673 & 0.359 & 0.569 & 0.440 & 0.690 \\
\midrule
\emph{Baseline: always small} & 0.774 & 0.000 & 0.000 & 0.000 & 0.500 \\
\bottomrule
\end{tabular}
\end{table}


For problem A every algorithm beats the baseline. The best, \AlgRegBest, reduces the MAE from \AlgRegBaseMAE to
\AlgRegBestMAE units, an improvement of \AlgRegGain\%. The spread between the top models is small, however, and
linear regression (\AlgLinMAE) is within a few units of the ensembles. Only the single decision tree
(\AlgDTMAE) clearly lags behind. For problem B every algorithm lies in a narrow band of ROC-AUC, from
\AlgClsWorstAUC (\AlgClsWorst) to \AlgClsBestAUC (\AlgClsBest). The algorithms differ mainly in where they place
the trade-off between precision and recall, which is a business choice made through the threshold, not a
difference in predictive power.

These results agree with \cref{ch:models}, which used a richer feature set and a different tuning procedure. In
both studies, (i) all machine-learning models beat the naive baselines, (ii) the leading regressors are
separated by only a few percent, so their exact ranking is not stable, (iii) linear models are competitive with
complex ensembles once good lag features exist, and (iv) large-order prediction is capped at a ROC-AUC of about
0.7 by the information in the data. Reaching the same conclusions with two independent implementations is a
form of triangulation that increases confidence in the findings \citep{shadish2002, kuhn2019}.

\begin{table}[htbp]\centering\small
\caption{Practical profile of the nine algorithms, based on the experiments}\label{tab:algoprofile}
\begin{tabularx}{\linewidth}{@{}lcccc>{\raggedright\arraybackslash}X@{}}
\toprule
\textbf{Algorithm} & \textbf{Type} & \textbf{Scaling} & \textbf{Explain.} & \textbf{Cost} & \textbf{Best use in this project} \\ \midrule
Linear regression & R & yes\textsuperscript{a} & high & very low & transparent demand baseline \\
Logistic regression & C & yes & high & very low & large-order probability, odds ratios \\
Decision tree & R/C & no & high & low & business rules, explaining drivers \\
Random forest & R/C & no & medium & medium & robust forecasts with ranges \\
Gradient boosting & R/C & no & medium & medium & accurate forecasts (tuned) \\
XGBoost & R/C & no & medium & low & accurate forecasts, early stopping \\
SVM / SVR & R/C & yes & low & high & linear SVR for demand; small data \\
KNN & R/C & yes & medium & high\textsuperscript{b} & case-based explanations \\
$K$-means & U & yes & high & low & product and market segments \\ \bottomrule
\end{tabularx}
\par\smallskip\footnotesize R = regression, C = classification, U = unsupervised; Explain.\ = explainability.
\textsuperscript{a}Not needed for the fit itself, but needed to compare coefficients.
\textsuperscript{b}At prediction time, because every prediction searches all stored rows.
\end{table}

\subsection{Lessons for practitioners}
\begin{enumerate}
  \item \textbf{Start with a baseline and a simple model.} A naive rule and a linear model set the bar. Here they
        were within a few percent of the best ensemble, at a fraction of the cost and with full explainability.
  \item \textbf{Tune on validation data, never on the test.} Every setting (depth, number of trees, learning
        rate, $C$, $\gamma$, $k$, threshold) was chosen on 2024. The test score is therefore an honest estimate.
  \item \textbf{Scale for distance-based methods.} SVM and KNN dropped to near-random performance without
        scaling. Tree methods do not need it.
  \item \textbf{Watch the gap between training and validation.} Deep trees, $k=1$ neighbours and boosting without
        early stopping all fit the training data almost perfectly and generalise worse.
  \item \textbf{Choose metrics that match the business.} With imbalanced classes, accuracy rewards a useless
        model. ROC-AUC, recall, precision and a deliberately chosen threshold are more meaningful.
  \item \textbf{When all algorithms agree, the limit is the data.} Seven classifiers with the same ROC-AUC show
        that better large-order prediction needs new information, such as customer history, not a different
        algorithm \citep{kelleher2020, geron2022}.
\end{enumerate}
\section{Decision-support insights}\label{sec:insights}
The final layer turns model outputs into recommendations. The module \texttt{insights.py} reads the metric
files and fills sentence templates with the current results. Every insight therefore stays up to date when new
data are processed, and none is typed by hand. The insights generated for this dataset are:

\begin{enumerate}[leftmargin=*,itemsep=4pt]
\item \textbf{Growth (What happened).} Net sales grew from 2020 to 2025 at 16.6\% a year, but units grew only 5.0\% a year; on average 11.3 percentage points of annual growth came from higher prices and mix, 5.0 from volume. \emph{Recommendation:} Track volume growth as a separate KPI; price-led growth is not sustainable indefinitely.
\item \textbf{Pricing (Why it happened).} After controlling for product, year and season, the estimated price elasticity of demand is +0.09 (95\% CI -0.28 to +0.46), i.e. not significantly different from zero. Discounts show no relationship with order size (Spearman rho = 0.001). \emph{Recommendation:} Review the discount budget (BDT 7.7 million over the period): discounts do not appear to buy volume.
\item \textbf{Seasonality (What happened).} December is the peak month (index 125, average month = 100) and February the weakest (86). The seasonal effect is statistically strong (Kruskal-Wallis p < 0.001). \emph{Recommendation:} Build stock and distributor capacity from October; schedule promotions and maintenance in February.
\item \textbf{Portfolio (What happened).} Muffin is the largest category (22.4\% of sales); the top 10 of 52 products deliver 38.5\% of revenue and 27 products form the 'A' class (80\% of sales). \emph{Recommendation:} Protect availability of A-class products first; review C-class products for rationalisation or repositioning.
\item \textbf{Portfolio (Segmentation).} Clustering finds 5 product segments; 'Premium \& gift packs' grows fastest (+7.8\% units a year, 33.3\% of sales). \emph{Recommendation:} Prioritise shelf space and marketing for 'Premium \& gift packs'.
\item \textbf{Market (What happened).} Dhaka alone generates 43.9\% of sales and Dhaka plus Chattogram 58.3\%. General Trade is the largest channel (45.6\%), while E-commerce is 5.9\%. \emph{Recommendation:} Reduce concentration risk by developing secondary cities and scaling E-commerce, which grows from a small base.
\item \textbf{Market (Segmentation).} The 'Emerging high-growth cells' segment (19 region-channel-customer cells, mainly Cumilla, Khulna, Mymensingh) is growing at +33.2\% a year, the fastest of the 4 market segments. \emph{Recommendation:} Allocate field-sales and digital-marketing effort to these emerging cells.
\item \textbf{Orders (Prediction).} Whether an order will be large ($\geq$ 10 units) is moderately predictable (best ROC-AUC 0.712, Logistic Regression); channel is by far the strongest driver: 48\% of distributor orders are large versus 8\% online. \emph{Recommendation:} Use the large-order score for delivery-vehicle and stock allocation, mainly at channel level.
\item \textbf{Forecast (What will happen).} The Support Vector Regression model forecasts next month's product demand with 11.1\% weighted error, 24\% lower RMSE than the seasonal-naive rule. Net sales for 2026-10 to 2027-09 are forecast at BDT 76.3 million (+12.1\% on the previous 12 months; 80\% range BDT 72.7 million - BDT 79.9 million). \emph{Recommendation:} Use the product-level forecast for the monthly production plan and the interval for safety stock.
\item \textbf{Risk (Anomalies).} 28 of 2,465 trading days were flagged as unusual. In a controlled test, rules plus Isolation Forest caught 91\% of 1,000 planted data errors. \emph{Recommendation:} Run the anomaly screen daily and investigate flagged days and orders before month-end reporting.
\end{enumerate}

Each insight is tagged by area and by the management question it answers. Each is paired with a concrete
action, so the dashboard serves as a \emph{decision-support system}, not merely a reporting tool.

\section{System evaluation}
The proposal asked for the system to be evaluated on predictive performance, usability, analytical usefulness,
visualisation quality and processing efficiency. \Cref{tab:syseval} summarises the evidence.

\begin{table}[htbp]\centering\small
\caption{Evaluation of the framework against the criteria in the proposal}\label{tab:syseval}
\begin{tabularx}{\linewidth}{@{}lX@{}}
\toprule
\textbf{Criterion} & \textbf{Evidence} \\ \midrule
Predictive performance & All ML forecasters beat naive baselines; best WAPE \BestRegWAPE\%; company-level MAPE \HWMAPE\%; classification ROC-AUC \ClsAUC (\cref{ch:models}); confirmed independently by the algorithm study (\cref{ch:algorithms}) \\
Analytical usefulness & Insights on growth quality, discount ineffectiveness, seasonality, concentration risk, segments and anomalies, each with a recommended action \\
Usability & One set of global filters; nine task-oriented pages with bookmarkable addresses; plain-language captions; downloadable forecasts; runs in any browser without installation \\
Visualisation quality & Principles of Few, Tufte and Knaflic: KPI cards, sorted bars, heat maps, intervals on forecasts, no 3-D or decorative charts, consistent palette \\
Processing efficiency & Data validation and cleaning of \NRows rows in under 20 seconds; warehouse load in \WarehouseSeconds s; full pipeline including model tuning in about 20 minutes on a standard laptop; dashboard data about 2 MB \\
Availability & Publicly deployed on free HTTPS platforms with automatic redeployment on every push (\cref{sec:deployment}) \\
Reproducibility & One command regenerates all data, models, figures, tables and numbers; fixed random seed; automated tests; version history on GitHub \\ \bottomrule
\end{tabularx}
\end{table}

A formal usability study with managers was beyond the time available. It is recommended as future work: a short
task-based test with a System Usability Scale questionnaire \citep{brooke1996} would complete the evaluation.

\section{Hypothesis testing results}\label{sec:hyptest}
\Cref{tab:hypdecision} brings together the decisions on the seven hypotheses of \cref{sec:hypotheses}, using
the tests defined in \cref{tab:hypmap}. Four of the seven alternative hypotheses (H1, H4, H6 and H7) are
supported and three (H2, H3 and H5) are not. The unsupported hypotheses are themselves among the most useful
findings for management, because they contradict three common assumptions: that customer type drives order
size, that discounts drive volume, and that higher prices reduce demand.

\begin{table}[htbp]\centering\small
\caption{Decisions on the research hypotheses}\label{tab:hypdecision}
\begin{tabularx}{\linewidth}{@{}lXll@{}}
\toprule
\textbf{Hyp.} & \textbf{Evidence} & \textbf{Effect} & \textbf{Decision} \\ \midrule
H1 & Order size differs by channel; $p<0.001$ & $\varepsilon^2=\EpsChannel$ & Supported \\
H2 & No difference by customer type; $p=\PCustType$ & $\varepsilon^2\approx 0$ & Not supported \\
H3 & Discount and order size uncorrelated; $\rho=\RhoDisc$, $p=\PDisc$ & $\rho^2\approx 0$ & Not supported \\
H4 & Detrended daily sales differ by month; $H=\KWMonthH$, $p<0.001$ & large & Supported \\
H5 & Elasticity \Elasticity, 95\% CI [\ElasticityLo, \ElasticityHi], $p=\ElasticityP$ & near zero & Not supported\textsuperscript{a} \\
H6 & ML error \GainVsSNaive\% below seasonal naive (WAPE \BestRegWAPE\% vs \SNaiveWAPE\%) & \GainVsSNaive\% & Supported \\
H7 & Large orders associated with channel; $p<0.001$ & $V=\CramerV$ & Supported \\ \bottomrule
\end{tabularx}
\par\smallskip\footnotesize \textsuperscript{a}Demand is price-\emph{inelastic} in the observed price range:
the elasticity is not distinguishable from zero.
\end{table}

The lack of support for H5 is consistent with the price-led growth seen in \cref{ch:eda}: prices rose by about
\CAGRPrice\% a year without a measurable loss of volume.

\section{Discussion: answers to the research questions}\label{ch:discussion}

\paragraph{RQ1: How can data analytics improve business decision-making?}
The framework shows improvements at all three levels of analytics. At the \emph{descriptive} level, it replaces
manual spreadsheet consolidation with validated, reconciled KPIs that update with one command. At the
\emph{diagnostic} level, it reveals relationships that simple reports would miss. Growth is mostly price- and
mix-led. Discounts are not associated with larger orders. Order size depends on the channel, not the customer
type. At the \emph{predictive} level, it provides product-level forecasts with quantified uncertainty. Each of
these changes a concrete decision: discount budgets, sales-force allocation, production plans and logistics
capacity.

\paragraph{RQ2: Which preprocessing techniques are most effective?}
Three techniques proved most valuable. The first is \emph{deriving} missing values from accounting identities,
which gives exact restoration and should always be tried before statistical imputation. The second is
\emph{robust, group-wise outlier detection} (a per-product IQR on the log scale), which respects skewed
distributions and different product scales. Combined with the policy of flagging rather than deleting unusual
values, it keeps genuine bulk orders. The third is \emph{leakage-safe feature engineering}, in which history
features are shifted before rolling and every learned parameter is fitted on training data only. Robust
methods were preferable to normal-theory methods throughout, because the business variables are highly skewed.

\paragraph{RQ3: Which algorithm performs best?}
No single algorithm is best for every problem, and the margins between good models are small. For one-month-ahead
product demand, \BestReg was most accurate. \BestRec was most robust in the multi-step setting. At the company
level, the classical Holt--Winters model beat all machine-learning alternatives. For classification, logistic
regression matched the tree ensembles. This pattern agrees with \citet{makridakis2018}. Good features, a sound
validation design and benchmarking matter more than the choice of algorithm. A combined approach, with
statistical totals reconciled to machine-learned product mixes, is the most useful in practice. The
algorithm-level study in \cref{ch:algorithms} reached the same conclusion with independent code and a simpler
feature set. There, \AlgRegBest had the lowest monthly error (MAE \AlgRegBestMAE units), but linear regression
(\AlgLinMAE) and the random forest (\AlgRFMAE) were within a few units of it, and the seven classifiers spanned
only \AlgClsWorstAUC to \AlgClsBestAUC in ROC-AUC. Because the exact ranking of the leading models changes with
the features and tuning, the choice among them should favour simplicity and explainability.

\paragraph{RQ4: What patterns exist in the historical data?}
Strong year-end seasonality peaks in \PeakMonth. Growth is steady but price-driven, and demand is
price-inelastic. Revenue is geographically concentrated in Dhaka and Chattogram, and the portfolio is
concentrated in about ten products. Order size differs strongly by channel. Discounts have no detectable volume
effect, and weekday effects are absent. Segmentation adds a portfolio view (premium and gift packs, core family
packs, value single-serves and niche lines) and a market view (core retail, bulk distribution, online and
regional, emerging high-growth).

\paragraph{RQ5: How accurately can the future be predicted?}
Monthly demand per product can be forecast with a weighted error of about \BestRegWAPE\% one month ahead and
about \BestRecWAPE\% up to 21 months ahead. Company-level monthly units can be forecast with an error of about
\HWMAPE\%. Whether an individual order will be large is only moderately predictable (ROC-AUC \ClsAUC). The
limit lies in the information available, not in the models.

\paragraph{RQ6: How should results be presented to managers?}
Through an interactive dashboard organised around the management questions. It should have global filters,
KPI cards and simple, sorted charts; it should always show uncertainty for forecasts; and it should state
insights in plain language, each with an action. Publishing the dashboard and a summary website on free
platforms removes cost and installation barriers. In this project both are live (\cref{tab:links}), every page
has its own shareable address, and continuous deployment means that managers always see the latest results
without any action on their part.

\paragraph{RQ7: How much does predictive analytics add to historical reporting?}
A historical report can at best repeat last year's pattern, which is the seasonal-naive rule. The
machine-learning forecast reduces its error by \GainVsSNaive\%. Put differently, a planner using last year's
figures would make errors about \SNaiveWAPE\% of volume in each product-month, compared with \BestRegWAPE\% for
the framework. Prediction intervals add information that historical reporting cannot provide at all: how much
safety stock is needed.

\section{Comparison with the literature}
\Cref{tab:litcompare} sets the main findings against the literature reviewed in \cref{ch:literature}. Most
findings confirm published evidence; two of them add context-specific knowledge about the Bangladeshi
packaged-cake market.

\begin{table}[htbp]\centering\small
\caption{Findings of this study compared with the literature}\label{tab:litcompare}
\begin{tabularx}{\linewidth}{@{}XXl@{}}
\toprule
\textbf{Finding of this study} & \textbf{Position in the literature} & \textbf{Relation} \\ \midrule
Holt--Winters beat all ML models for company-level monthly units & Simple statistical methods are hard to beat on aggregate series \citep{makridakis2018, hyndman2021} & Confirms \\
ML models beat naive baselines for product-level demand, with small margins between them & Feature design and validation matter more than the algorithm \citep{kuhn2019, kuhn2013} & Confirms \\
Random validation would have overstated accuracy & Time-ordered validation is required for temporal data \citep{bergmeir2012, kaufman2012} & Confirms \\
Logistic regression matched tree ensembles for large orders & Linear models remain competitive when signal is limited \citep{hastie2009, james2021} & Confirms \\
Several metrics were needed to compare classifiers under imbalance & Single metrics hide classifier differences \citep{japkowicz2011} & Confirms \\
Segments are interpretable but only moderately separated (silhouette $\approx$ \ProdSil) & Real market data rarely form well-separated clusters \citep{kaufman1990, kotler2016} & Confirms \\
Discounts have no measurable effect on order size; demand is price-inelastic & Context-specific; requires experimental confirmation \citep{wooldridge2020} & Extends \\
Order size depends on channel, not customer type & Context-specific evidence for the Bangladeshi FMCG market & Extends \\ \bottomrule
\end{tabularx}
\end{table}

\section{Chapter summary}
Growth was price-led, seasonality strong and sales concentrated, while discounts had no measurable effect.
Product demand was forecast with a WAPE of \BestRegWAPE\% (\GainVsSNaive\% better than the seasonal-naive rule)
and company units with a Holt--Winters MAPE of \HWMAPE\%. The algorithm-level study confirmed these results
independently; four of seven hypotheses were supported and all research questions answered.
